\section{Généralité sur les courbes elliptiques}

\subsection{Rappels et notations}

On suppose connues les notions de théorie des courbes elliptiques présentent dans le cours du master. On énonce ici les définitions et notations utilisées dans la suite du rapport. 


\begin{Def}
	On défini une courbe elliptique $E$ à isomorphisme prêt. Par abus de notation, on désigne par $E$ un ensemble de courbes elliptiques de même $j$-invariant. Un modèle de courbe elliptique est une courbe elliptique explicite donnée par une équation. Pour préciser le choix d'un modèle de courbe de $E$, on le donne par son équation. 
\end{Def}

\begin{Rem}
	On considère toujours des modèles de courbes elliptiques donnés par des équation de Weierstrass réduites $E : y^2 = x^3 + ax + b$. On notera $\OR$ le point neutre, ou $\OR_{E}$ s'il y a besoin de précision (pour différentier des ordres par exemple). 
\end{Rem}

\begin{Def}
	Soit $E : y^2 = x^3 + ax + b$ un modèle de courbe elliptique. Il est défini sur $\K$ si $a$, $b\in\K$. On notera $E(\K)$ l'ensemble des points du modèles définis sur $\K$, et $E(\Kb)$ l'ensemble des points du modèles définis sur $\Kb$. Si $P\in(\Kb)$ est un points du modèle $E$, on note $x(P)$ son abscisse et $y(P)$ son ordonnée. Une courbe elliptique $E$ est dites définies sur un corps $\K$ s'il existe au moins un modèle de courbes elliptique de $E$ dont l'équation est définie sur $\K$.
\end{Def}


\begin{Def}
	Soit $E$ et $E'$ deux courbes elliptiques. Une isogénie $\phi : E\rightarrow E'$ est défini à isomorphisme prêt. Par abus de notation, on désigne par $\phi : E\rightarrow E'$ un ensemble d'isogénie égales à isomorphisme prêt au départ et à l'arrivée. Une isogénie $\phi$ est sous forme explicite lorsque l'on spécifie un modèle de départ et d'arrivé. $\phi$ est un endomorphisme lorsque $E = E'$. 
\end{Def}

\begin{Rem}
	Ce choix de notation sera pratique pour les énoncés théoriques, mais lorsque l'on voudra calculer une isogénie, On voudra plus précisément calculer des formes explicite d'isogénie.  
\end{Rem}

\begin{Ex}
	\begin{itemize}
		\item Soit $E : y^2 = x^3 + ax + b$ un modèle de courbe elliptique. Si $P\in E$, on désigne par $\tau_{P}$ la translation par $P$.
		\item Soit $E$ une courbe elliptique. On note $[m]$ l'endomorphisme de $E$ consistant en la multiplication scalaire des points par $m$. On note $E[m]$ l'ensemble des points de $m$-torsions de $E$ (la notation $E[m]$ est donc aussi définie à isomorphisme prêt si l'on ne précise pas de modèle.)
	\end{itemize}
\end{Ex}

\begin{Def} 
	Soit $E$ une courbe elliptique définie sur $\K$. On défini l'ensemble des fonctions rationnelles sur $E$, noté $\Kb(E)$, à isomorphisme de $E$ prêt. Un élément de $\Kb(E)$ est donc une fonction rationnelle $f : E\rightarrow \Kb$ définie à isomorphisme de $E$ prêt.
\end{Def}

\begin{Def}
	Soit $E$, $E'$ deux courbes elliptiques définies sur un corps $\K$. Soit $\phi : E\rightarrow E'$ est une isogénie explicite. on défini son application contravariante (ou Pullback) par
	\begin{equation*}
		\phi^{*} : \Kb(E') \rightarrow \Kb(E), f\mapsto f\circ\phi
	\end{equation*} 
	On désigne par $\hat{\phi} : E'\rightarrow E$ son isogénie duale.
\end{Def}

\begin{Def}
	Soit $E$, $E'$ deux courbes elliptiques définies sur un corps $\K$. Soit $\phi : E\rightarrow E'$ une isogénie. Une forme explicite de $\phi$ est dite définie sur $\K$ si les modèles de départ et d'arrivé sont définis sur $\K$ et que les formules sous formes rationnelles sont aussi définies sur $\K$. L'isogénie $\phi$ est définie sur $\K$ si elle admet une forme explicite définie sur $\K$. 
\end{Def}

\begin{Rem}
	Si une isogénie $\phi : E\rightarrow E'$ est défini sur $\K$ et que l'on se donne une forme explicite définie sur $\K$, le noyau de $\phi$ n'est pas forcément défini sur $\K$.
\end{Rem}











\subsection{Noyaux d'isogénie}


Lorsque deux courbes elliptiques $E$ et $E'$ sont défini sur un corps $\K$, on sait que l'on peut associé à une isogénie $\phi : E\mapsto E'$ une extension $\Kb(E)/\phi^{*}\left( \Kb(E')\right)$. Par exemple les notions de degré ou de séparabilité de $\phi$ sont définis par rapport à cette extension. On commence ici par appliquer la théorie de Galois à cette extension. On va observer que le groupe de Galois $Gal\left( \Kb(E)/\phi^{*}\left( \Kb(E')\right)\right)$ est isomorphes au noyaux de $\phi$.


\begin{The}\cite{Sil09}[III 4.10]
	
	Soit $E$, $E'$ deux courbes elliptiques défini sur un corps $\K$, et $\phi : E\rightarrow E'$ une isogénie non nulle. 
	\begin{enumerate}
		\item L'application 
		\begin{equation*}
			Ker(\phi) \rightarrow Aut\left( \left[ \Kb(E) : \phi^{*}\left( \Kb(E')\right) \right] \right) , P\mapsto \tau_{P}^{*}
		\end{equation*} 
		est un isomorphisme.
		\item Si $\phi$ est séparable, l'extension $\Kb(E)/\phi^{*}\left( \Kb(E')\right)$ est galoisienne, i.e $deg(\phi) = |Ker(\phi)|$
	\end{enumerate}
\end{The}

\begin{Coro}\cite{Sil09}[III 4.12]
	
	Soit $E$ une courbe elliptique et $G$ un sous groupe fini de $E$. Il existe une unique courbe elliptique $E'$ et une isogénie séparable $\phi : E \rightarrow E'$ telles que $Ker(\phi) = G$.
\end{Coro}

\begin{Rem}
	On fera attention à l'unicité dans ce corollaire. En pratique le modèle de départ sera donné, et il faudra fixer un modèle d'arrivé, et expliciter un moyen de calculer l'image d'un point $P$ du modèle de départ. 
\end{Rem} 

\begin{Def}
	Soit $E$ une courbe elliptique et $G$ un sous groupe fini de $E$. L'unique courbe elliptique image d'une isogénie de noyau $G$ est noté $E/G$.
\end{Def}

Lorsque le noyau $G$ est donné comme sous groupe d'un modèle de courbe elliptique, on peut calculer une version explicite de l'isogénie de noyau $G$ avec une formule dû à Vélu donnant un modèle d'arrivé et l'image d'un point $P$:

\begin{Prop}{Formule de Vélu}\cite{Velu71}
	
	Soit $E : y^2 = x^3 + ax + b$ un modèle de courbe elliptique défini sur un corps $\K$. Soit $G\subset E(\Kb)$ un sous-groupe fini. L'isogénie $\phi : E\rightarrow E/G$ de noyaux $G$ admet une forme explicite partant du modèle $E$, où l'image d'un point $P\in E(\Kb)$ est définie par :
	
	\begin{equation*}
		\phi(P) = \left( x(P) + \sum_{Q\in G\setminus{\OR}} x(P+Q) - x(Q)\;,\;y(P) + \sum_{Q\in G\setminus{\OR}} y(P+Q) - y(Q)\right) 
	\end{equation*} 
	
	Et le modèle d'arrivé est donné par l'équation $E' : y^2 = x^3 + a'x + b'$ avec
	\begin{equation*}
		a' = a - 5\sum_{Q\in G\setminus{\OR}}(3 (x(Q))^2 + a )
	\end{equation*}
	
	\begin{equation*}
		b' = b - 7\sum_{Q\in G\setminus{\OR}}(5 (x(Q))^3 + 3ax(Q) + b )
	\end{equation*}
\end{Prop}

Lorsqu'une courbe elliptique est définie sur un corps $\K$, on souhaitera calculer des formes explicites d'isogénies définies sur $\K$. On énonce un critère sur $G$ pour assurer que l'isogénie $\phi : E\rightarrow E/G$ est définie sur $\K$.

\begin{Def}
	Soit $E : y^2 = x^3 + ax + b$ un modèle de courbe elliptique définie sur un corps $\K$. Soit $G$ un sous groupe fini de $E(\Kb)$, et $P\in E(\Kb)$. Le groupe $Gal(\Kb/\K)$ agit sur $P$ et $G$ de la façon suivante
	\begin{align*}
		\forall\sigma\in Gal(\Kb/\K),\;& P^{\sigma} = \left( \sigma(x(P)), \sigma(y(P))\right)\\
		& G^{\sigma} = \left\lbrace P^{\sigma},\; P\in G \right\rbrace 
	\end{align*}
	
	
\end{Def}

\begin{Prop}\cite{Sil09}[II 4.12, 4.13.2]
	
	Soit $E$ et $E'$ deux courbes elliptiques définies sur un corps $\K$. Une isogénie $\phi : E\rightarrow E'$ est définie sur $\K$ si et seulement si son noyau est stable sous l'action de $Gal(\Kb/\K)$, c'est à dire
	\begin{equation*}
		\forall\sigma\in Gal(\Kb/\K),\; Ker(\phi)^{\sigma} \subset Ker(\phi)
	\end{equation*}
\end{Prop}

On s'intéresse maintenant aux isogénies séparables de degré donné, c'est à dire dont on connaît le cardinal du noyau.

\begin{Def}
	On dira qu'une $\phi : E\rightarrow E'$ entre deux courbe elliptique est une $l$-isogénie si elle est de degré $l$. 
\end{Def}

\begin{Prop}
	Si $\phi : E\rightarrow E'$ est une $l$-isogénie entre les courbes elliptiques $E$ et $E'$, alors $Ker(\phi)\subset E\left[ l\right] $
\end{Prop}

\begin{proof}
	On rappelle la propriété fondamentale de l'isogénie duale : $\left[ l\right] = \hat{\phi}\circ\phi$. Ainsi si $P$ est un point d'un modèle de $E$, et si $P\in Ker(\phi)$, alors $lP = \hat{\phi}\circ\phi(P) = \hat{\phi}(\OR) = \OR $. Donc $P$ est un point de $l$-torsion de son modèle.
\end{proof}

On en déduit le nombre exacte de $l$-isogénie séparable sur une courbe elliptique $E$ :

\begin{Prop}
	Soit $E$ une courbe elliptique et $l\in\N$.
	\begin{enumerate}
		\item(rappel) $E[l] \simeq \Z/l\Z\times\Z/l\Z$
		\item Il existe exactement $l+1$ sous groupes d'ordre $l$ dans $\Z/l\Z\times\Z/l\Z$.
		\item Il existe exactement $l+1$ $l$-isogénie séparable sur $E$.
	\end{enumerate}
\end{Prop}


\begin{proof}
	Soit $P$, $Q$ une base de $E\left[ l\right]$ comme $\Z/l\Z$ module. $P$ et $Q$ sont ordre $l$ et sont libre sur $\Z/l\Z$, donc $\forall i\in\left[ 0 , l - 1\right] $, $P + iQ$ est d'ordre $l$. Ces points engendres des groupes disjoints (à l'exception du neutre $\OR$). On a donc trouvé $l + 1$ sous groupes d'ordre $l$ qui n'ont deux à deux que le neutre en commun. Ils sont constitués au total de $1 + (l+1)*(l - 1) = l^2$ points distincts. Comme $E[l]$ est de cardinal $l^2$, on a trouvé exactement tout les sous groupes d'ordre $l$. Chacun d'eux défini exactement une isogénie séparable sur $E$.
\end{proof}











\subsection{Invariant différentiel et normalisation d'isogénie}

Étant donné un modèle d'une courbe elliptique $E$ et un sous groupe $G$ fini, on souhaite définir un modèle d'arrivé pour une forme explicite de l'isogénie de noyaux $G$, de façon normalisé pour tout $E$, $G$. Pour cela on va définir l'invariant différentiel d'un modèle de courbe elliptique. Le modèle d'arrivé sera déterminé par son invariant différentiel. 

On commence par une théorie générale des invariants différentiels sur une courbe algébrique. Dans toute la suite, $C$ désigne une courbe algébrique projective supposée lisse, et $\Kb(C)$ l'ensemble de fonction rationnel sur $C$. 

\begin{Def}
	L'ensemble des formes différentielles de $C$, noté $\Omega_{C}$, est un $Kb(C)$-espace vectoriel engendré par des symboles $dx$ où $x$ parcours $\Kb(C)$.
	$dx$ est une dérivée formelle : on demande que trois relations soit vérifiée
	\begin{enumerate}
		\item $d(x+y) = dx + dy$ pour tout $x, y \in\Kb(C)$
		\item $d(xy) = xdy + ydx$ pour tout $x, y \in\Kb(C)$
		\item $da = 0$ pour tout $a\in\Kb$
	\end{enumerate}
\end{Def}

\begin{Prop}\cite{Sil09}[II 4.2]
	
	\begin{enumerate}
		\item $\Omega_{C}$ est de dimension 1 sur $\Kb(C)$
		\item Pour $x \in \Kb(C)$, $dx$ est une base de $\Omega_{C}$ si et seulement si l'extension $\Kb(C) / \Kb(x)$ est fini et séparable. 
	\end{enumerate}
\end{Prop}


\begin{Rem} On sait que l'extension $\Kb(C) / \Kb(x)$ est fini (théorème de normalisation de Noether) dès que $x\in\Kb(C)$ est non constant. En caractéristique 0 $dx$ est donc toujours une base de $Omega_{C}$ car la séparabilité de l'extension est garantie. 
\end{Rem}

On peut étendre l'application contravariante d'un morphisme de courbe algébrique $\phi$ à aux formes différentielles. Cela donne notamment un critère de séparabilité qui nous sera utile par la suite.

\begin{Def}
	Soit $\phi : C \rightarrow C'$ un morphisme entre courbes algébriques $C$ et $C'$. On défini $\phi^{*} : \Omega_{C'} \rightarrow \Omega_{C}$ de la façon suivante :
	\begin{equation*}
		\phi^{*}\left(\sum f_idx_i \right) = \sum \left( \phi^{*}(f_i) \right) d\left( \phi^{*}(x_i)\right) 
	\end{equation*}
\end{Def}

\begin{Prop}\cite{Sil09}[II 4.2]
	
	Soit $\phi : C \rightarrow C'$ un morphisme entre courbes algébriques $C$ et $C'$. $\phi$ est séparable si et seulement si $\phi^{*} : \Omega_{C'} \rightarrow \Omega_{C}$ est injective, ce qui équivaut à ce que $\phi^{*}$ soit non nulle. 
\end{Prop}

Sur une courbe algébrique lisse $C$, $f\in\Kb(C)$ vérifie $div(f) = O$ si et seulement si $f$ est constante. On va généraliser la notion de diviseur aux formes différentielles. Certaines vont alors avoir un diviseur nulle : ceux seront les invariants différentiels. 

\begin{Prop}\cite{Sil09}[II 1.4]
	
	Soit $C$ une courbe algébrique lisse. Soit $P\in C$ et $u\in \Kb(C)$ une uniformisante en $P$. Alors $du$ est une base de $\Omega_{C}$. 
\end{Prop}

\begin{Def}
	Soit $C$ une courbe algébrique lisse et $\omega\in\Omega_{C}$. Soit $P\in C$ et $u\in \Kb(C)$ une uniformisante en $P$. Il existe $g\in\Kb(C)$ tel que $\omega = gdu$. On note $\frac{\omega}{du} = g$
\end{Def}

\begin{Prop}\cite{Sil09}[III 4.3]
	
	Soit $C$ une courbe algébrique lisse et $\omega\in\Omega_{C}$. Soit $P\in C$ et $u\in \Kb(C)$ une uniformisante en $P$. L'ordre en $P$ de $\frac{\omega}{du} = g$ ne dépend pas du choix d'uniformisante $u$. On note donc $ord_{P}(\omega) = ord_{P}(\frac{\omega}{du})$. Alors $ord_{P}(\omega)$ est nul sauf pour un nombre fini de points $P$.
\end{Prop}



\begin{Def}
	Avec les notation précédentes, On défini $div(\omega) = \sum_{P\in C} ord_{P}(\omega)(P) \in Div(C)$. On dira que $\omega\in\Omega_{C}$ est holomorphe si $div(\omega)\geq 0$, et sans zéros si $div(\omega)\leq 0$. Un invariant différentiel $\omega$ est une forme différentielle holomorphe sans zéros, c'est à dire telle que $div(\omega) = 0$.
\end{Def}

\begin{Rem}
	
	\begin{itemize}
		\item On peut montrer que l'application $div$ se comporte sur $\Omega_{C}$ comme sur $\Kb(C)$, en particulier si $f\in\Kb(C)$ et $\omega\in\Omega_{C}$ alors $div(f\omega) = div(f) + div(\omega)$. Cela implique que tout les diviseurs de formes différentielles sont égaux dans le groupe $Pic(C)$. On appelle cette classe la classe canonique de $Pic(C)$. 
		
		\item S'il existe un invariant différentiel, alors la classe canonique est la classe trivial. On peut montrer avec le théorème de Riemann-Roch que s'il existe un invariant différentiel sur $C$, alors $C$ est de genre $1$.
		
	\end{itemize}
\end{Rem}

L'existence d'un invariant différentiel est en réalité une propriété propre au genre $1$.

\begin{Prop}\cite{Sil09}[II 4.6]
	
	Soit $E : y^2 = x^3 + ax + b$ un modèle courbe elliptique. Alors $\omega = \frac{dx}{2y}\in\Omega_{E}$ est un invariant différentiel de $E$.
\end{Prop}

\begin{Def}
	Soit $E : y^2 = x^3 + ax + b$ un modèle courbe elliptique. On note $\omega_{E}$ l'invariant différentiel du modèle, défini par $\omega_{E} = \frac{dx}{2y}$. 
\end{Def}

\begin{Rem}
	Puisque si $f\in\Kb(C)$ et $\omega\in\Omega_{C}$ alors $div(f\omega) = div(f) + div(\omega)$, tout les invariants différentiels d'une courbe elliptique $E$ sont égaux à une constante multiplicative prêt.
\end{Rem}

Étant donné un modèle $E$ de courbe elliptique d'invariant différentiel $\omega_{E}$, il est facile de trouver un modèle isomorphe $E'$ d'invariant différentiel $u\omega_{E}$ pour tout $u\in\Kb$

\begin{Prop}\cite{Sil09}[III table 3.1]
	
	Soit $E : y^2 = x^3 + ax + b$ un modèle courbe elliptique définie sur un corps $\K$, et $u\in\Kb$. Le modèle $E' : y^2 = x^3 + a'x + b'$ obtenue à partir de $E$ par le changement de variable $(x,y)\mapsto(u^{2}x,u^{3}y)$ est isomorphe à $E$ et vérifie $\omega_{E'} = u\omega_{E}$.
\end{Prop}

On va se servir de la proposition 2.12 pour définir une convention de normalisation des formes explicite d'isogénie.

\begin{Def}
	Soit $\phi : E\rightarrow E'$ une isogénie sous forme explicite, $\omega_{E}$ et $\omega_{E'}$ les invariants différentiels associés respectivement aux modèles de départ et d'arrivé. On dira que cette forme explicite de $\phi$ est normalisée si $\phi^{*}(w_{E'}) = w_E$
	
\end{Def}

On énonce alors une version explicite du corollaire 2.2.

\begin{Prop}
	Soit $E : y^2 = x^3 + ax + b$ un modèle courbe elliptique et $G$ un sous-groupe fini de $E$. Soit $\phi : E\rightarrow E/G$ l'isogénie de noyau $G$. La courbe elliptique $E/G$ admet un unique modèle $E' : y^2 = x^3 + a'x + b'$ tel que la forme explicite de $\phi$ de départ $E$ et d'arrivé $E'$ soit normalisée. 
\end{Prop}

\begin{Rem}
	La formule de Vélu énoncée ne fournit pas une forme explicite normalisée.
\end{Rem}


L'action de l'application contravariante d'une isogénie sur un invariant différentiel $\omega_{E}$ va s'avérer très simple, ce qui sera utile par la suite pour étudier les endomorphismes. 


\begin{The}\cite{Sil09}[III, 5.2]
	
	Soit $E$, $E'$ deux courbes elliptiques définie sur un corps $\K$, et $\phi$, $\psi : E \rightarrow E'$ deux isogénie. Soit $\omega$ un invariant différentiel d'un modèle de $E'$. Alors
	\begin{equation*}
		\left( \phi + \psi\right) ^{*}\left( \omega\right) = \phi^{*}\left( \omega\right) + \psi^{*}\left( \omega\right)
	\end{equation*} 
\end{The}


\begin{Coro}\cite{Sil09}[III, 5.3]
	
	Soit $E$ une courbes elliptiques sur $\K$, $\omega$ un invariant différentiel d'un modèle de $E$. Alors pour $m\in\Z^{*}$,
	$[m]^{*}(\omega) = m\omega$. En particulier $[m]^{*}(\omega) \neq 0 $ et on retrouve que $[m]$ est séparable.
\end{Coro}


\begin{Coro}\cite{Sil09}[III, 5.6]
	
	Soit $E$ une courbe elliptique sur $\K$, $\omega$ un invariant différentiel d'un modèle de $E$. Pour tout endomorphisme $\phi\in End(E)$, soit $a_{\phi}\in\Kb(E)$ tel que $\phi^{*}(\omega) = a_{\phi}\omega$. On considère
	
	\begin{equation*}
		f : End(E)\rightarrow\Kb(E) , \phi\longmapsto a_{\phi}
	\end{equation*}
	Alors
	\begin{enumerate}
		\item $a_{\phi}\in\Kb$ et $f: End(E)\rightarrow\Kb$ est bien définie. 
		\item $f$ est un morphisme d'anneaux
		\item $Ker(f)$ est exactement l'ensemble des endomorphisme inséparable de $E$
		\item Si $car(K) = 0$ alors $End(E)$ est commutatif
	\end{enumerate}
	
\end{Coro}



\subsection{Anneaux d'endomorphisme}

\begin{Def}
	Soit $E$ une courbe elliptique définie sur un corps $\K$. Tout les modèles de $E$ ont des anneaux d'endomorphismes isomorphes, et l'on note $End(E)$ l'anneau d'endomorphisme de $E$
\end{Def}

\begin{Rem}
	Si $E : y^2 = x^3 + ax + b$ et $E' : y^2 = x^3 + a'x + b'$ sont deux modèles isomorphes, une isogénie ayant une forme explicite $\phi : E\rightarrow E'$ est un endomorphisme, même si les modèles $E$ et $E'$ sont différents. Il sera souvent nécessaire pour normaliser d'écrire un endomorphisme sous une forme explicite où les modèles de départ et d'arrivé sont distincts. 
\end{Rem}

Il n'existe que $3$ structure possible pour un anneau d'endomorphisme de courbe elliptique.

\begin{The}\cite{Sil09}[III 9.4]
	
	Soit $E$ une courbe elliptique définie sur un corps $\K$. L’anneau $End(E)$ est isomorphe soit à $\Z$, soit à un ordre d'un corps quadratique imaginaire, soit à un ordre dans une algèbre de quaternions.
\end{The}

\begin{Rem} Soit $E$ une courbe elliptique définie sur un corps $\K$.
	\begin{itemize}
		\item Si $\K$ est de caractéristique nulle, $End(E)$ est un anneaux commutatif, donc soit isomorphe à $\Z$, soit à un ordre d'un corps quadratique imaginaire.
		\item Si $\K$ est de caractéristique $p>0$, l'endomorphisme de Frobenius est un morphisme qui ne correspond pas à une multiplication scalaire, donc $End(E)$ n'est pas isomorphe à $\Z$.
	\end{itemize}
\end{Rem}

Ainsi la seule structure d'anneaux d'endomorphisme qui peut exister en toute caractéristique et celle d'ordre de corps quadratique imaginaire.  

\begin{Def}
	Une courbe elliptique $E$ est dites "à multiplication complexe", ou "ordinaire", si $End(E)\sim \OR_{\Delta}$ où $\OR_{\Delta}$ est un ordre quadratique imaginaire de discriminant $\Delta$. 
\end{Def}

\begin{Def}
	Une courbe elliptique $E$ définie sur un corps $\K$ de caractéristique non nulle est supersingulière si $End(E)$ est isomorphe à un ordre dans une algèbre de quaternion.  
\end{Def}